\section{Quadratic self-twists are excluded for the descended type-II points}
\label{sec:cubic-typeii-noncm}

\begin{lemma}[AC-K29-nonCM-trace: a two-character trace with quadratic symmetry]
\label{AC-K29-nonCM-trace}
Let $k$ be a field of characteristic different from $2$, let
$G$ be a group, and let $a,b:G\to k^\times$ be characters.
Suppose a nontrivial quadratic character $\epsilon:G\to\{\pm1\}$
satisfies
\[
       a(g)+b(g)=\epsilon(g)\bigl(a(g)+b(g)\bigr)
                         \qquad(g\in G).                 \tag{NC.1}
\]
Then $a/b=\epsilon$.
\end{lemma}
\begin{proof}
For $g\notin\ker\epsilon$, (NC.1) and $2\ne0$ give
$a(g)/b(g)=-1$.  Fix one such $g_0$.  For every
$h\in\ker\epsilon$, both $g_0$ and $g_0h$ are outside
the kernel.  Hence
$(a/b)(h)=(a/b)(g_0h)/(a/b)(g_0)=1$.
These two cosets exhaust $G$, proving the assertion.
No semisimplicity of a representation with this trace is required.
\end{proof}

\begin{theorem}[AC-K29-nonCM: no quadratic self-twist at the actual type-II points]
\label{AC-K29-nonCM}
Retain the literal old datum $D$, all its \texttt{GD} hypotheses,
and the actual cubic-to-old point transport in
AC-K29-kill, AC-K29-point, and AC-G48-cubic.
In particular, the coefficient prime satisfies
$\ell\notin\{2,3,269\}$, $e\geq1$, and the point has
the actual stable ordinary flag used in AC-G48.
Let
\[
       V:G_{\mathbb Q}\longrightarrow\mathrm{GL}_2(E)
                                                               \tag{NC.2}
\]
be its normalized descended endpoint representation, after a
finite coefficient extension if necessary.  Its determinant
is $\eta_K$, and its strict residual semisimplification is
\[
 \overline V^{\mathrm{ss}}
       =(\overline\eta_K\,\overline\gamma)
                        \oplus\overline\gamma^{-1},
 \qquad \gamma=\kappa_\ell^e\psi.                       \tag{NC.3}
\]
Then $V$ has no nontrivial quadratic self-twist over
$\overline{\mathbb Q}_\ell$.
In particular it is not induced from a character of any
quadratic number field.  The same conclusions hold after
any scalar character twist of $V$, including the source
normalization $V\otimes\kappa_\ell^{-e}$.

Consequently the newform supplied by the retained pointwise
modularity input in AC-K29-typeII has no CM quadratic
self-twist.  This supplies the non-CM prerequisite for an
Euler-system theorem; it supplies neither such a theorem's
other hypotheses nor a vanishing conclusion.
\end{theorem}
\begin{proof}
\proofstep{1. The actual residual and local data.}
In the source normalization, if $B(g)$ is the determinant-one
endpoint block of the old semidirect representation, then
\[
  V(g)=B(g)\operatorname{diag}(-1,1)^{\varepsilon_K(g)},
       \qquad \eta_K(g)=(-1)^{\varepsilon_K(g)}.           \tag{NC.4}
\]
This is the explicit \texttt{Wfun} of
\path{UnitaryThreeGroups/ExistsLieElementMiddleBlockPart.lean}
and the \texttt{W\_facts}/\texttt{typeII\_facts} proof in
\path{UnitaryThreeGroups/CoboundaryLemmaNoColumn.lean}.
Reducing (NC.4), whose old endpoint seed is
$\operatorname{diag}(\overline\gamma,
\overline\gamma^{-1})$, proves (NC.3), including the
$\eta_K$ in the first character.

The source \texttt{GD.l\_cop} makes $K$ unramified at
$\ell$.  The fields \texttt{GD.ine} and \texttt{GD.gen}
in \path{HeightAndVanishing/Counts.lean}, lines 15 and 22,
give, on the inertia group $I_\ell$,
\[
  \gamma^2|_{I_\ell}=\kappa_\ell^{2e}|_{I_\ell},
  \quad
  \overline\gamma^2|_{I_\ell}
       =\overline\kappa_\ell^{2e}|_{I_\ell}\ne1.          \tag{NC.5}
\]
Here \texttt{ine} says that the $\ell$-primary conductor
component of $\psi$ has square $1$; the prime-to-$\ell$
conductor component is unramified at $\ell$.  The last
inequality uses the surjectivity of the mod-$\ell$
cyclotomic character on inertia and $\ell-1\nmid2e$.

The stable ordinary line in AC-G48-TypeII and its strict
upper-triangular basis give two continuous local characters
$\lambda_1,\lambda_2:G_{\mathbb Q_\ell}\to E^\times$
with
\[
 \operatorname{tr}V|_{G_{\mathbb Q_\ell}}
       =\lambda_1+\lambda_2,
 \qquad
 (\lambda_1/\lambda_2)|_{I_\ell}
                          =\kappa_\ell^{2e}|_{I_\ell}.    \tag{NC.6}
\]
Indeed, (G48.8) writes the diagonal as
$\eta_K\delta,\delta^{-1}$ with
$\delta|_{I_\ell}=\gamma|_{I_\ell}$, and $\eta_K$ is
unramified there.  This uses the actual stable flag, not
only its semisimplified inertia eigencharacters.  The
ratio in (NC.6) has infinite image: the cyclotomic inertia
image contains $1+\ell\mathbb Z_\ell$, and a nonzero
integer power remains infinite.  Equivalently the two
local characters have distinct Hodge--Tate weights
$e$ and $-e$; the argument below needs only the displayed
inertia identity.

\proofstep{2. A hypothetical global self-twist is ramified at $\ell$.}
Suppose $V\simeq V\otimes\epsilon$ for a nontrivial
quadratic character of $G_{\mathbb Q}$.  Equality of traces
holds for every group element.  Reduce that equality in
the original integral lattice for $V$.  No integral
intertwiner is needed: both sides are already equal
integral traces.  Since $\ell\ne2$, the reduction of
$\epsilon$ remains a nontrivial quadratic character.
Apply the preceding lemma to (NC.3).  It gives
\[
           \overline\epsilon
              =\overline\eta_K\,\overline\gamma^2.      \tag{NC.7}
\]
By (NC.5) and the unramifiedness of $\eta_K$, this character
is nontrivial on $I_\ell$.  Thus $\epsilon$ itself is
nontrivial on $G_{\mathbb Q_\ell}$, in fact on inertia.

\proofstep{3. The local flag contradicts that self-twist.}
Restrict the same trace equality to $G_{\mathbb Q_\ell}$.
The two local characters in (NC.6) and the preceding lemma,
now in characteristic zero, imply
\[
                   \lambda_1/\lambda_2
                         =\epsilon|_{G_{\mathbb Q_\ell}}.
                                                               \tag{NC.8}
\]
The right side has order at most two, whereas the left
side has infinite inertia image by (NC.6).  This is the
required contradiction.  No choice of CM field, local
splitting assertion, or residual lattice comparison enters
this step.

\proofstep{4. Induction and the modular normalization.}
If $V=\operatorname{Ind}_{G_F}^{G_{\mathbb Q}}\theta$
for a quadratic extension $F/\mathbb Q$, write the induced
space in its two cosets.  The matrix
$\operatorname{diag}(1,-1)$ commutes with the diagonal
action of $G_F$ and changes the sign of the action of the
other coset.  It is an explicit intertwiner between $V$
and $V\otimes\eta_F$.  This contradicts the result just
proved.  A scalar twist cancels from a putative self-twist
isomorphism, so the conclusion is unchanged under every
scalar twist.  In the customary modular formulation, a
CM quadratic self-twist likewise gives the displayed
quadratic trace symmetry of its attached representation.
Thus the source's modular normalization has no such CM
self-twist.  All uses of source automorphy in that final
interpretation retain the hypotheses of AC-K29-typeII.
\end{proof}

\paragraph{Scope for the remaining cubic gate.}
This theorem closes a representation-theoretic prerequisite
for every type-II point satisfying the established point
transport hypotheses.  It does not construct the cubic
Type I automorphic seed, remove the projected mixed tangent
condition in AC-K29-tangent, or treat the excluded coefficient
primes.  In particular, it is not a substitute for an
Euler-system or Bloch--Kato theorem in the actual weight and
character range.

\paragraph{Geometric dictionary.}
The representation $V$ is an $\ell$-adic local system on
the relevant arithmetic open.  A quadratic character is
a rank-one local system associated with a quadratic finite
\'etale cover, and induction from that cover is its
rank-two direct image.  The theorem excludes the resulting
quadratic tensor symmetry for the actual local systems
under consideration.  For finite reductions this language
uses the already audited finite-cover dictionary of
SGA~1, Expos\'e~V, \S\S4 and~7.  The exclusion itself
comes from the two-coset trace calculation and the actual
ordinary inertia characters above; it is not an additional
result attributed to that categorical dictionary.
